% !PRINTED 475
\weberpage{475}

%<<CONT
was wieder die identische Permutation ist. \gesperrt{Es ist also die zu
$\pi_a^{-1}$ entgegengesetzte Permutation wieder die ursprüngliche
Permutation $\pi_a$.}

Ist $\pi_a \pi_b = \pi_c$, so ist
\[
\pi_c^{-1} = \pi_b^{-1} \pi_a^{-1},
\]
denn es ist $\pi_a \pi_b \pi_b^{-1} \pi_a^{-1}$ die identische
Permutation.

Bezeichnen wir die identische Permutation mit $\pi_0$, so
ist also

\weq{10}{\pi_a \pi_a^{-1} = \pi_a^{-1} \pi_a = \pi_0,}

und $\pi_0$ ist sich selbst entgegengesetzt (wodurch nicht
ausgeschlossen ist, dass es noch andere sich selbst entgegengesetzte
Permutationen giebt). Durch die Composition mit der identischen
Permutation wird keine andere Permutation geändert, also

\weq{11}{\pi_0 \pi_a = \pi_a \pi_0 = \pi_a.}

\gesperrt{2. Ist $\pi_c$ aus $\pi_a$ und $\pi_b$ zusammengesetzt, also}

\weq{12}{\pi_a \pi_b = \pi_c,}

\gesperrt{so ist nicht nur $\pi_c$ aus $\pi_a$ und $\pi_b$, sondern
auch $\pi_a$ aus $\pi_b$ und $\pi_c$, und $\pi_b$ aus $\pi_a$ und
$\pi_c$ völlig und eindeutig bestimmt.}

Denn es folgt aus (12)

\weq{13}{\pi_a = \pi_c \pi_b^{-1}, \quad \pi_b = \pi_a^{-1} \pi_c.}

\gesperrt{3. Ein aus der Gesammtheit aller $\Pi(m)$ Permutationen von
$m$ Elementen herausgegriffenes System
\[
Q = \pi, \pi', \pi'', \pi''' \ldots,
\]
das so in sich abgeschlossen und vollendet ist, dass je zwei und
folglich auch beliebig viele Permutationen von $Q$ durch
Zusammensetzung immer wieder eine Permutation desselben Systems
ergeben, heisst eine Permutationsgruppe. Unter dem Grad einer
Permutationsgruppe verstehen wir die Anzahl der Permutationen, die sie
enthält.}

Die Gesammtheit aller $\Pi(m)$ Permutationen bildet eine solche
Gruppe. Ebenso ist die einzige identische Permutation eine Gruppe für
sich. Was zwischen diesen beiden extremen Fällen von Gruppen noch
liegen kann, das wird uns in den folgenden Abschnitten beschäftigen.
